\emph{Fungsi}~$f\colon A\to B$ masangaké saben anggota
\emph{ranah}~$A$ karo persis siji anggota \emph{kodomain}~$B$.
Yèn $x\in A$, $y$ kang dipasangaké~$f$ marang~$x$ diarani
\emph{nilai} $f(x)$ ing \emph{argumèn}~$x$. Yèn $A$ minangka
himpunan pasangan, fungsi~$f$ bisa dianggep nampa rong argumèn.
\emph{Jangkauan} $\ran{f}$ iku subhimpunan saka~$B$ kang
ngemot kabèh nilai~$f$.

Yèn $\ran{f}=B$, $f$ diarani \emph{surjèktif}. Saben
argumèn~$x$ dipasangaké marang siji nilai~$f(x)$. Yèn ora ana
rong argumèn béda kang dipasangaké marang nilai kang padha,
$f$ diarani \emph{injèktif}. Fungsi kang injèktif lan surjèktif
diarani \emph{bijèktif}.

Fungsi bijèktif nduwèni siji \emph{fungsi invers} $f^{-1}$.
Fungsi bisa diranté. Fungsi $g\circ f$ iku \emph{komposisi}
$f$ banjur~$g$. Komposisi fungsi-fungsi injèktif iku injèktif,
lan komposisi fungsi-fungsi surjèktif iku surjèktif.
$f^{-1}\circ f$ iku fungsi identitas.

Yèn $f$ kena ora nduwèni nilai kanggo sawatara $x\in A$,
fungsi iku diarani \emph{fungsi parsial}. Kanggo fungsi
parsial $f\colon A\pto B$, kita nulis $f(x)\fdefined$ yèn
$f(x)$ \emph{katetepaké}, yaiku yèn $f$ nduwèni nilai ing
argumèn~$x$. Yèn ora, $f(x)$ \emph{ora katetepaké},
$f(x)\fundefined$. Saben fungsi~$f$, kalebu fungsi parsial,
nduwèni \emph{graf}~$R_f$, yaiku relasi kang dumadi yèn lan
mung yèn $f(x)=y$.
